% GENERATED FILE. Edit canonical lessons, then run npm run generate:books.
% canonical-source-hash: fnv1a64:38d1c5869cf19a9c
% canonical-lessons: RU-C05-brat, RU-C05-sprashivat, RU-C05-pomogat, RU-C05-lyubit

\chapter{Taking, Asking, Helping, Loving}
\label{ch:russian-verbs-of-the-hand}
\hlchaptermodality{5}

\begin{chapteropening}
\textbf{By the end of this chapter:} \emph{I can say I take, ask, help and love in Russian, put any verb I own after \ru{люблю} as a bare infinitive, and tell a suppletive aspect pair like \ru{брать}·\ru{взять} from an ordinary one like \ru{спрашивать}·\ru{спросить}.}

Say \ru{я} \ru{люблю} with its inserted \ru{л}, then run every verb you own through it — \ru{я} \ru{люблю} \ru{читать}, \ru{писать}, \ru{помогать}, \ru{жить} — and close by naming the chapter's four partners and picking out \ru{взять} as the one built from a different root.
\end{chapteropening}

\section[\texorpdfstring{\ru{брать} (brat')}{брать (brat')}]{\ru{брать} — \textquotedblleft{}to take\textquotedblright{}}

\label{lesson:RU-C05-brat}



\begin{quote}
A short verb, an English cousin you carry every day, and a partner that shares not one letter with it.
\end{quote}

\subsection*{You'll want to know first — \ru{брать}}
\begin{quote}
\textbf{\ru{брать}} — \emph{brat'} — \textbf{to take}
\end{quote}

Five letters, all of them yours, and the familiar \textbf{-\ru{ть}}.

\begin{quote}
\emph{ya berú} — \textbf{I take.}
\end{quote}

The stem grows a vowel: \textbf{\ru{бр}-} in the infinitive, \textbf{\ru{бер}-} once an ending arrives. Then \emph{ty beryósh}, so this verb sits in the \emph{znat'} family, and that \textbf{-u} is the plain form of the \textbf{-\ru{ю}} you have been saying.

Three sentences, all ordinary Russian:

\begin{quote}
\emph{Ya chitáyu. Ya pishú. Ya berú.}
\end{quote}

\begin{grammarlens}[title={Grammar Lens: a pair with nothing in common}]
The partners so far were recognisable: \emph{chitát' · prochitát'} added a prefix, \emph{ponimát' · \ru{понять}} reshaped a middle. Now the hard case.

\textbf{\ru{Брать}}'s finished partner is \emph{vzyat'} — \emph{ya vaz'mú}. Not a prefixed \emph{\ru{брать}}, not a squeezed one: a \textbf{different word entirely}, drafted in as half of the pair.

Grammarians call this \textbf{suppletion}, and you have met it once: \emph{\ru{идти}} takes \emph{shol} for its past, from another root, exactly as English \emph{go} takes \textbf{went}. Here the same trick builds an aspect pair instead of a past tense.

There is a quiet reward in it. \emph{Vzyat'} is \emph{vz-} on the very root that built \emph{ponimát'}, the ancient root meaning \emph{take}. So \emph{\ru{брать}}'s finished partner is family with the verb for \emph{understand}, and \emph{\ru{брать}} is the outsider.
\end{grammarlens}

\begin{cousinweb}
\textbf{\ru{брать}} is Indo-European *\emph{b\textsuperscript{h}er-}, \textquotedblleft{}\textbf{to carry, to bear}\textquotedblright{}, one of the richest roots English owns: \textbf{bear}, \textbf{birth}, \textbf{burden}, \textbf{bier}; Latin \emph{ferre} behind \textbf{transfer}, \textbf{refer}, \textbf{fertile}; Greek \emph{phérein} behind \textbf{metaphor}, a carrying-across. Russian keeps the load too — \emph{brémya} is a burden, and \emph{berémennaya}, \emph{pregnant}, is literally \emph{carrying}.

Now the trap, because this track flags both kinds. \textbf{\ru{Брат}} means \textbf{brother}, sits one soft sign from \emph{\ru{брать}}, and is \textbf{not related}: \emph{\ru{брат}} is *\emph{b\textsuperscript{h}réh\textsubscript{2}ter}, English \emph{brother}. And \emph{byt'} is *\emph{b\textsuperscript{h}uH-}, English \emph{be}. Three roots that all begin \textbf{b\textsuperscript{h}} — a family accent, not a shared word.
\end{cousinweb}

\subsection*{Your turn}
Say these aloud:

\begin{itemize}
\raggedright
  \item \textquotedblleft{}\ru{брать}\textquotedblright{} then \textquotedblleft{}ya berú\textquotedblright{} — hear the stem grow its vowel
  \item \textquotedblleft{}Ya chitáyu. Ya pishú. Ya berú.\textquotedblright{} — \emph{ya chitáyu, ya pishú, ya berú}
  \item the odd pair — \textquotedblleft{}\ru{брать} … vzyat'\textquotedblright{}, two words sharing nothing
  \item the other suppletion — \textquotedblleft{}\ru{идти} … shol\textquotedblright{}, \textquotedblleft{}go … went\textquotedblright{}
  \item \textquotedblleft{}\ru{брать}\textquotedblright{} then English \textquotedblleft{}bear, birth, burden, transfer, metaphor\textquotedblright{}
\end{itemize}

\begin{tcolorbox}[breakable,colback=teal!4,colframe=teal!35!black,title={Before you move on}]
Say \textquotedblleft{}I take.\textquotedblright{} (\emph{Ya berú}.) Name \emph{\ru{брать}}'s finished partner and say what is strange about it, then give the English pair that behaves the same way. (\emph{Vzyat'}, a different root; \emph{go} and \textbf{went}, as \emph{\ru{идти}} takes \emph{shol}.) Which verb of yours is \emph{vzyat'} related to? (\emph{Ponimát'} — both on the old root for \emph{take}.) And is \emph{\ru{брать}} related to \emph{\ru{брат}}? (\textbf{No} — different roots, and \emph{byt'} is a third.)
\end{tcolorbox}
\section[\texorpdfstring{\ru{спрашивать} (spráshivat')}{спрашивать (spráshivat')}]{\ru{спрашивать} — \textquotedblleft{}to ask\textquotedblright{}}

\label{lesson:RU-C05-sprashivat}



\begin{quote}
You have asked a question in Russian since Chapter 2 without owning the verb for asking. Here it is.
\end{quote}

\subsection*{You'll want to know first — \ru{спрашивать}}
\begin{quote}
\textbf{\ru{спрашивать}} — \emph{spráshivat'} — \textbf{to ask (a question)}
\end{quote}

Long on the page, easy in the mouth: the stress lands on \textbf{\ru{спра}-} and stays, \textbf{SPRA-shi-vat'}. The \textbf{\ru{ш}} in the middle is the letter \emph{pishú} introduced.

\begin{quote}
\emph{ya spráshivayu} — \textbf{I ask.}
\end{quote}

The same \textbf{-\ru{ю}}, and \emph{ty spráshivayesh} puts it in the \emph{znat'} family. Set it beside your other verb for opening your mouth: \textbf{\ru{говорить}} is to speak, \textbf{\ru{спрашивать}} is to speak with a question mark on the end. And one distinction English hides: Russian asks a \textbf{question} with \emph{\ru{спрашивать}}, but \textbf{for a thing} with \emph{\ru{просить}}.

\begin{grammarlens}[title={Grammar Lens: an ordinary pair, after a very strange one}]
\emph{\ru{Брать}} took \emph{vzyat'}, a partner from a different root. That was the worst case, not the norm. \textbf{\ru{Спрашивать}}'s partner is \textbf{\ru{спросить}} — \emph{ya sproshú, ty sprósish} — one root worn into two shapes, the ongoing member long, the finished short.
\end{grammarlens}

\begin{culture}
The moment you ask, Russian makes you choose. \emph{Kak vas zovút?} takes \textbf{\ru{вас}} for a stranger and \textbf{\ru{тебя}} for a friend — a choice that sits inside \emph{ty} and \emph{vy}, the split English lost when \textbf{thou}, \emph{ty}'s cousin, left ordinary use. Other languages make the same move by other grammar: French \textbf{vous} and Russian \emph{vy} are the plural doing polite duty, German \textbf{Sie} borrows the third person, Spanish \textbf{usted} grew out of a title.

And recall what that question literally asks: \textbf{how} they call you, not what your name is — an action where English reaches for a possession, as French and Spanish also do.
\end{culture}

\begin{cousinweb}
Strip the prefix and \textbf{\ru{спрашивать}} rests on the root of \textbf{\ru{просить}}, \textquotedblleft{}to request\textquotedblright{}: Indo-European *\emph{pre\'{k}-}, \textquotedblleft{}\textbf{to ask}.\textquotedblright{} In Latin it became \emph{precārī}, \textquotedblleft{}to pray\textquotedblright{}, and English took the lot: \textbf{pray}, \textbf{prayer}, \textbf{precarious}, \textbf{deprecate}. German \textbf{fragen} is the Germanic outcome, so \emph{\ru{спрашивать}} and \emph{fragen} are one inherited word.

English \textbf{ask} is not in this family; the cousin kept here is \emph{pray}.
\end{cousinweb}

\subsection*{Your turn}
Say these aloud:

\begin{itemize}
\raggedright
  \item \textquotedblleft{}\ru{спрашивать}\textquotedblright{} then \textquotedblleft{}\ru{я} \ru{спрашиваю}\textquotedblright{} — and hear the \emph{\ru{ш}} of \emph{pishú}
  \item the two mouths — \textquotedblleft{}\ru{я} \ru{говорю} … \ru{я} \ru{спрашиваю}\textquotedblright{}
  \item \textquotedblleft{}Kak vas zovút\textquotedblright{} then \textquotedblleft{}Kak tebyá zovút\textquotedblright{} — and choose which fits
\end{itemize}

\begin{tcolorbox}[breakable,colback=teal!4,colframe=teal!35!black,title={Before you move on}]
Say \textquotedblleft{}I ask.\textquotedblright{} (\emph{Ya spráshivayu}.) Name the partner and say how this pair differs from \emph{\ru{брать} · vzyat'}. (\textbf{\ru{Спросить}} — one root in two shapes.) In \emph{Kak vas zovút?}, what decides between \textbf{\ru{вас}} and \textbf{\ru{тебя}}? (Whether the person is \emph{vy} or \emph{ty} — and \emph{ty} is English \emph{thou}.) And give an English cousin of *\emph{pre\'{k}-}. (\emph{Pray}, \emph{prayer}, \emph{precarious}.)
\end{tcolorbox}
\section[\texorpdfstring{\ru{помогать} (pomogát')}{помогать (pomogát')}]{\ru{помогать} — \textquotedblleft{}to help\textquotedblright{}}

\label{lesson:RU-C05-pomogat}



\begin{quote}
This verb is built out of the word for \emph{being able}, and that word is English \textbf{may}. Helping, in Russian, is lending someone your can.
\end{quote}

\subsection*{You'll want to know first — \ru{помогать}}
\begin{quote}
\textbf{\ru{помогать}} — \emph{pomogát'} — \textbf{to help}
\end{quote}

The \textbf{\ru{г}} is the letter \emph{\ru{говорить}} introduced: Greek gamma, a hard \textbf{g} as in \emph{go}. Both \textbf{\ru{о}}s lean towards \emph{a} away from the stress, as \emph{\ru{говорить}}'s do — \textbf{pa-ma-GAT'}.

\begin{quote}
\emph{ya pamagáyu} — \textbf{I help.}
\end{quote}

Same \textbf{-\ru{ю}}, and \emph{ty pamagáyesh} files it with \emph{znat'}.

One honest limit: the person you help takes a case this book has not reached, so keep the sentence bare. \emph{Ya pamagáyu} alone is ordinary Russian.

\begin{grammarlens}[title={Grammar Lens: a partner that does not end in -\ru{ть}}]
Every infinitive so far ended in \textbf{-\ru{ть}}. The partner of \emph{\ru{помогать}} does not: it is \emph{pomóch'}, ending in \textbf{-ch'}, and \emph{ya pamagú, ty pamózhesh}.

A small handful of Russian verbs end that way, and they are built on old roots about power. The shape is a signpost, not an exception to memorise blind.

Three pairs cover the chapter so far: \textbf{\ru{брать} · vzyat'}, \textbf{spráshivat' · \ru{спросить}}, \textbf{\ru{помогать}} · \emph{pomóch'}.
\end{grammarlens}

\begin{cousinweb}
\textbf{\ru{помогать}} is \textbf{\ru{по}-} plus \textbf{-\ru{могать}}, and under it sits \emph{moch'}, \textquotedblleft{}to be able\textquotedblright{} — \emph{ya magú}, \textquotedblleft{}I can.\textquotedblright{}

That root is Indo-European *\emph{mag\textsuperscript{h}-}, \textquotedblleft{}\textbf{to be able, to have power}\textquotedblright{}, and English holds several pieces of it:

\begin{itemize}
\raggedright
  \item \textbf{may} and \textbf{might} — the ordinary modal, and the noun for strength in \emph{with might and main}.
  \item \textbf{dismay} — to strip someone of their power to act.
  \item German \textbf{mögen} and \textbf{Macht}; and Greek \emph{mēkhan\'{\={e}}}, \textquotedblleft{}a contrivance\textquotedblright{}, which is English \textbf{machine} and \textbf{mechanic}.
\end{itemize}

Note what did \textbf{not} happen here. \emph{Magú} and \emph{may} still mean the same thing after five thousand years. That is not the usual outcome: \emph{zhit'}'s root handed English \textbf{quick}, which once meant \textbf{alive} and now means \emph{fast}. Some roots keep their sense, some drift out of recognition, and \emph{spráshivat'}'s went from \emph{ask} all the way to \emph{pray}. Each one has to be checked rather than guessed.
\end{cousinweb}

\subsection*{Your turn}
Say these aloud:

\begin{itemize}
\raggedright
  \item \textquotedblleft{}\ru{помогать}\textquotedblright{} — \emph{pa-ma-GAT'}, both \emph{o}s leaning to \emph{a}
  \item \textquotedblleft{}Ya pamagáyu\textquotedblright{}
  \item the three pairs — \textquotedblleft{}\ru{брать}·vzyat', spráshivat'·\ru{спросить}, \ru{помогать}·pomóch'\textquotedblright{}
  \item \textquotedblleft{}magú\textquotedblright{} then English \textquotedblleft{}may, might\textquotedblright{} — one meaning, two languages
\end{itemize}

\begin{tcolorbox}[breakable,colback=teal!4,colframe=teal!35!black,title={Before you move on}]
Say \textquotedblleft{}I help.\textquotedblright{} (\emph{Ya pamagáyu}.) Which verb is hiding inside it, and what does it mean? (\emph{Moch'} — to be able; \emph{ya magú}.) Name the partner and say what is odd about its ending. (\emph{Pomóch'}, in \textbf{-ch'} rather than \textbf{-\ru{ть}}.) Give two English cousins of *\emph{mag\textsuperscript{h}-}. (\emph{May}, \emph{might}, \emph{dismay}, \emph{machine}.) And which English word from your verbs no longer means what its root meant? (\textbf{Quick}, from \emph{zhit'}'s root — it meant \emph{alive}.)
\end{tcolorbox}
\section[\texorpdfstring{\ru{любить} (lyubít')}{любить (lyubít')}]{\ru{любить} — \textquotedblleft{}to love, and to like\textquotedblright{}}

\label{lesson:RU-C05-lyubit}



\begin{quote}
The last verb of the chapter is English \textbf{love}, and it lets you say every verb you own in a real sentence.
\end{quote}

\subsection*{You'll want to know first — \ru{любить}}
\begin{quote}
\textbf{\ru{любить}} — \emph{lyubít'} — \textbf{to love}, and \textbf{to like}
\end{quote}

One verb for both: with a person, \emph{love}; with an activity, closer to \emph{like}.

\begin{quote}
\emph{ya lyublyú} — \textbf{I love.}
\end{quote}

And \emph{ty lyúbish} — an \textbf{-ish}, so \emph{\ru{любить}} joins \emph{\ru{говорить}} and \emph{\ru{видеть}}, not \emph{znat'}. The stress moves too: \emph{lyublyú} on the ending, \emph{lyúbish} on the stem.

\begin{grammarlens}[title={Grammar Lens: the letter that appears out of nowhere}]
\emph{Lyubít'} has no \textbf{\ru{л}} after the \textbf{\ru{б}}. \emph{Lyublyú} does.

You have met this shape. \textbf{\ru{Видеть}} swapped \textbf{\ru{д}} for \emph{zh} in \emph{vízhu} and nowhere else. Here it happens by insertion: stems ending in \textbf{\ru{б}, \ru{п}, \ru{в}} or the letters for \emph{m} and \emph{f} push an \textbf{\ru{л}} in — but \textbf{only in the \emph{I} form}.

So \textbf{\ru{люблю}}, then \emph{lyúbish}, \textbf{\ru{любят}}, with no \textbf{\ru{л}} at all. Learn the \emph{I} form and the rest behaves.
\end{grammarlens}

\begin{grammarlens}[title={Grammar Lens: now say everything at once}]
\emph{Lyubít'} takes a bare infinitive, as English \emph{like} does — the dictionary form goes straight in:

\begin{quote}
\emph{Ya lyublyú chitát'.} — I like reading. \emph{Ya lyublyú pisát'.} — I like writing. \emph{Ya lyublyú dúmat'.} — I like thinking. \emph{Ya lyublyú pomogát'.} — I like helping. \emph{Ya lyublyú zhit'.} — I love living.
\end{quote}
\end{grammarlens}

\begin{cousinweb}
\textbf{\ru{любить}} is Indo-European *\emph{leub\textsuperscript{h}-}, and English \textbf{love} is the same word — no borrowing, no detour. The family is wider than it looks: \textbf{lief}; \textbf{believe} (\emph{be-} plus the root — to hold a thing dear enough to trust); \textbf{leave} in \emph{by your leave}; and Latin \emph{libet}, \textquotedblleft{}it pleases\textquotedblright{}, behind \textbf{libido}.

Its partner \textbf{\ru{полюбить}} deserves an honest line: not \emph{to finish loving} but \textbf{to come to love}, the moment the feeling starts. A pair does not always mean \emph{doing} against \emph{done} — which is why they get a later chapter of their own.

And the root thread: \textbf{\ru{бр}-} carries (\emph{bear}), \textbf{-\ru{прос}-} asks (\emph{pray}), \textbf{mog-} is able (\emph{might}), \textbf{\ru{люб}-} holds dear (\emph{love}) — with \textbf{\ru{пис}-} from last chapter still scratching its mark (\emph{picture}).
\end{cousinweb}

\subsection*{Your turn}
Say these aloud:

\begin{itemize}
\raggedright
  \item \textquotedblleft{}\ru{любить}\textquotedblright{} then \textquotedblleft{}\ru{я} \ru{люблю}\textquotedblright{}, then \textquotedblleft{}ty lyúbish\textquotedblright{} without the \emph{\ru{л}}
  \item \textquotedblleft{}Ya lyublyú chitát'. Ya lyublyú pisát'. Ya lyublyú pomogát'. Ya lyublyú zhit'.\textquotedblright{}
  \item the four pairs — \textquotedblleft{}\ru{брать}·vzyat', spráshivat'·\ru{спросить}, pomogát'·pomóch', \ru{любить}·\ru{полюбить}\textquotedblright{}
\end{itemize}

\begin{tcolorbox}[breakable,colback=teal!4,colframe=teal!35!black,title={Before you move on}]
Say \textquotedblleft{}I love\textquotedblright{}, then \textquotedblleft{}I like reading.\textquotedblright{} (\emph{Ya lyublyú}; \emph{Ya lyublyú chitát'}.) Where is the \textbf{\ru{л}} in \emph{ty lyúbish}, and which earlier verb behaved this way? (Gone — the insert hits only the \emph{I} form, as \textbf{\ru{д} $\to$ zh} in \emph{vízhu}.) Name the chapter's four partners, and which is a different root from its verb. (\emph{Vzyat'}, \textbf{\ru{спросить}}, \emph{pomóch'}, \textbf{\ru{полюбить}}; \emph{vzyat'}.) And give the English cousins of \emph{\ru{любить}}, \emph{\ru{брать}}, \emph{spráshivat'}, \emph{pomogát'} and \emph{\ru{писать}}. (\textbf{Love}, \emph{bear}, \emph{pray}, \emph{might}, \emph{picture}.)
\end{tcolorbox}
